\documentclass[b5j,fleqn]{jarticle}
\usepackage{zennyakuUniv}
\begin{document}
\body{
\Title{newcrown3 p50}
\para

\sent{Subtitles or Dubbing?}
{字幕か、それとも吹き替えか。}

\para

\sent{I prefer dubbing when I watch movies made in foreign languages.}
{私は外国語で作られた映画を見るとき、吹き替えのほうを好みます。}

\sent{First, with dubbed movies, you can focus on the visuals, not small words at the bottom of the screen.}
{第一に、吹き替え映画なら、画面の下にある小さな文字ではなく、映像に集中できます。}

\sent{You will catch the details of the movie.}
{映画の細かい部分まで捉えることができるでしょう。}

\sent{Second, it is easier to understand the story and characters.}
{第二に、物語や登場人物をより理解しやすくなります。}

\sent{Voice actors catch the original characters' tone of voice, feelings, and emotions.}
{声優は、元の登場人物の声の調子や気持ち、感情を捉えます。}

\sent{So you can concentrate on the story.}
{そのため、物語に集中できます。}

\sent{Dubbing is the better way to enjoy a movie.}
{映画を楽しむには、吹き替えのほうがよい方法です。}
}
\end{document}
